\section{Tag encodings with regional zero resets}
\label{sec:peps_hole_encoder}

Let $V$ be a finite set of sites, with finite coordinate sets $A_v$, and
choose a configuration $z\in\prod_{v\in V}A_v$. Write
$\mathcal H_R=\C^{\prod_{v\in R}A_v}$ and
$\ket{0_R}=\bigotimes_{v\in R}\ket{z_v}$.
For nested regions $R_0\subseteq\cdots\subseteq R_{n-1}$ and subspaces
$S_j\subseteq\mathcal H_{R_j}$, let $\widetilde S_j$ be the projected images
of Definition~\ref{def:peps_nested_cylinder_images}, and put
$\mathcal E=\sum_{j<n}\mathcal C_{R_j}(S_j)$.
The tag and zero-reset construction is the linear map appearing in the
\href{https://github.com/openai/math/blob/adc7f1241b42e322a6451854ab7e4b4c146bf78a/preprints/Polynomial-PEPS-approximation-of-gapped-square-grid-ground-states-September-24-2026/build/sections/05-frames.tex\#L50-L69}{adaptive hole encoding}
of the September 24, 2026 polynomial-PEPS manuscript.
% Source: 05-frames.tex, lines 13--97, at
% adc7f1241b42e322a6451854ab7e4b4c146bf78a; eq:hole-projector,
% eq:hole-encoder, eq:encoder-contraction, and the empty-outer convention.
% These entries concern one hole's finite-dimensional encoder only.
% They do not assert the geometric approximation input, ownership changes,
% the multi-hole frame estimate, or the full source Lemma 6.3.

\begin{definition}[Coordinate reset]
    \label{def:peps_coordinate_reset}
    \lean{TNLean.PEPS.coordinateResetMatrix}
    \leanok
    For a finite coordinate set $A$, a coordinate $a_0\in A$, and an
    arbitrary vector $v\in\C^A$, define
    $Z_{a_0}(v)=\ket{a_0}\bra v$.
    Its entry at $(a,b)$ is $\overline{v(b)}$ if $a=a_0$, and zero otherwise.
\end{definition}

\begin{theorem}[Reset Gram operators and basis projectors]
    \label{thm:peps_coordinate_reset_gram}
    \lean{
        TNLean.PEPS.coordinateResetMatrix_gram,
        TNLean.PEPS.coordinateRangeProjector_eq_sum_basis
    }
    \leanok
    \uses{def:peps_coordinate_reset}
    For any $v\in\C^A$ and any orthonormal basis $(b_k)_{k\in B}$ of a
    subspace $S\subseteq\C^A$,
    \begin{align}
        Z_{a_0}(v)^\dagger Z_{a_0}(v) & =\ketbra v v,
        \label{eq:peps_reset_gram}
    \end{align}
    and
    \begin{align}
        P_S & =\sum_{k\in B}\ketbra{b_k}{b_k}.
        \label{eq:peps_inside_basis_projector}
    \end{align}
    The reset identity does not require $v$ to be normalized.
\end{theorem}
\begin{proof}
    \leanok
    The only nonzero output coordinate of $Z_{a_0}(v)$ is $a_0$.
    Its contribution to the Gram entry at $(a,b)$ is
    $v(a)\overline{v(b)}$. Expansion in an orthonormal basis gives
    $P_S\psi=\sum_k b_k\langle b_k,\psi\rangle$.
\end{proof}

\begin{definition}[Tag register and zero-reset encoding]
    \label{def:peps_hole_encoder}
    \lean{
        TNLean.PEPS.HoleEncoderTag,
        TNLean.PEPS.holeEncoderBasis,
        TNLean.PEPS.holeEncoder,
        TNLean.PEPS.holeEncoder_apply,
        TNLean.PEPS.holeEncoder_eq_zero_of_not_reset,
        TNLean.PEPS.holeEncoder_eq_zero_of_complement_ne
    }
    \leanok
    \uses{def:peps_nested_cylinder_images,def:peps_coordinate_reset}
    Set $d_j=\dim\widetilde S_j$ and choose an orthonormal basis
    $(v_{j\ell})_{0\leq\ell<d_j}$ of each $\widetilde S_j$.
    Define the tag set $T=\{(j,\ell):0\leq j<n,\ 0\leq\ell<d_j\}$
    and the tag space $\mathcal J=\C^T$. The encoding
    $K\colon\mathcal H_V\to\mathcal J\otimes\mathcal H_V$ is
    \begin{align}
        K & =\sum_{j<n}\sum_{\ell<d_j}\ket{j,\ell}\otimes
        \left[(\ket{0_{R_j}}\bra{v_{j\ell}})_{R_j}
            \otimes I_{V\setminus R_j}\right].
        \label{eq:peps_hole_encoder}
    \end{align}
    In particular, the output retains one raw register at every site.
    For input and output configurations $y,x\in\prod_{v\in V}A_v$,
    its matrix coefficient is
    \begin{align}
        K_{((j,\ell),x),y}
        \label{eq:peps_hole_encoder_coefficient}
         & =
        \begin{cases}
            \overline{v_{j\ell}(y\vert_{R_j})},
               & x\vert_{R_j}=z\vert_{R_j},\quad
            x\vert_{V\setminus R_j}=y\vert_{V\setminus R_j}, \\
            0, & \text{otherwise}.
        \end{cases}
    \end{align}
    The reset region in this formula is $R_j$, determined by the tag.
    At a fixed tag $t=(j,\ell)$, the following tensor product depicts its summand.
    % Formula: |j,l> tensor (|0_Rj><v_jl| tensor I_(V\Rj)).
    % Ink-to-index: the top row is the inside covector and reset ket;
    % the middle row is the tag ket; the bottom box is the complement identity.
    % No legs are contracted internally. The left ports are y|Rj and y|Rc;
    % the right ports are x|Rj, (j,l), and x|Rc, all physical coordinates.
    % Boundary signature: zero virtual ports and five physical ports,
    % comprising two input and three output bundles. The two output raw
    % bundles regroup to one full configuration x. The tag sum is external.
    \begin{center}
        \tenkzkernel
        \begin{tenkz}[rows={wire,wire,wire}, cols=2, bonds=none]
            \tn[at=(1,1), name=holeInsideBra, skin=box,
                ports={180:physical:$y\vert_{R_j}$}]{\bra{v_{j\ell}}}
            \tn[at=3 e of holeInsideBra, name=holeResetKet, skin=box,
                ports={0:physical:$x\vert_{R_j}$}]{\ket{0_{R_j}}}
            \tn[at=2 s of holeResetKet, name=holeTagKet, skin=box,
                ports={0:physical:$t$}]{\ket{j,\ell}}
            \tn[at=4 s of holeInsideBra, name=holeComplement, skin=box,
                ports={180:physical:$y\vert_{V\setminus R_j}$,
                        0:physical:$x\vert_{V\setminus R_j}$}]{I_{V\setminus R_j}}
        \end{tenkz}
    \end{center}
\end{definition}

\begin{theorem}[Dimension of the tag register]
    \label{thm:peps_hole_encoder_tag_dimension}
    \lean{
        TNLean.PEPS.finrank_coordinateSubspaceES,
        TNLean.PEPS.card_holeEncoderTag,
        TNLean.PEPS.card_holeEncoderTag_le
    }
    \leanok
    \uses{def:peps_hole_encoder}
    Identifying coordinate vectors with their Euclidean Hilbert space
    preserves the dimension of every inside subspace.
    The tag dimension satisfies
    \begin{align}
        \dim\mathcal J=|T|
         & =\sum_{j<n}\dim\widetilde S_j
        \leq\sum_{j<n}\dim S_j.
        \label{eq:peps_hole_encoder_tag_dimension}
    \end{align}
\end{theorem}
\begin{proof}
    \leanok
    \uses{thm:peps_nested_cylinder_orthogonalization}
    The tags form a disjoint union of sets of cardinality $d_j$.
    Sum the inside-dimension inequalities
    $\dim\widetilde S_j\leq\dim S_j$.
\end{proof}

\begin{theorem}[The Gram operator of the encoding]
    \label{thm:peps_hole_encoder_gram}
    \lean{TNLean.PEPS.holeEncoder_gram}
    \leanok
    \uses{def:peps_hole_encoder}
    The encoding has Gram operator
    \begin{align}
        K^\dagger K
         & =\sum_{j<n}P_{\widetilde S_j}\otimes I_{V\setminus R_j}
        =P_{\mathcal E}.
        \label{eq:peps_hole_encoder_gram}
    \end{align}
\end{theorem}
\begin{proof}
    \leanok
    \uses{thm:peps_coordinate_reset_gram,thm:peps_nested_cylinder_projector_sum}
    Orthogonality of the tag vectors removes all products between distinct
    summands of~(\ref{eq:peps_hole_encoder}). For each remaining summand,
    $\braket{0_{R_j}}{0_{R_j}}=1$, so its Gram operator is
    $\ketbra{v_{j\ell}}{v_{j\ell}}\otimes I_{V\setminus R_j}$.
    Summing over the orthonormal basis of $\widetilde S_j$ gives
    $P_{\widetilde S_j}\otimes I_{V\setminus R_j}$.
    Theorem~\ref{thm:peps_nested_cylinder_projector_sum} identifies the sum
    over $j$ with $P_{\mathcal E}$.
\end{proof}

\begin{theorem}[Contraction and projected norm]
    \label{thm:peps_hole_encoder_norm}
    \lean{
        TNLean.PEPS.holeEncoder_opNorm_le_one,
        TNLean.PEPS.holeEncoder_norm_le,
        TNLean.PEPS.holeEncoder_norm_eq_projector
    }
    \leanok
    \uses{def:peps_hole_encoder}
    For every $\psi\in\mathcal H_V$,
    \begin{align}
        \|K\psi\| & =\|P_{\mathcal E}\psi\|\leq\|\psi\|.
        \label{eq:peps_hole_encoder_norm}
    \end{align}
    Thus $\|K\|\leq1$.
\end{theorem}
\begin{proof}
    \leanok
    \uses{thm:peps_hole_encoder_gram}
    Since $P_{\mathcal E}$ is an orthogonal projector, the $C^*$-identity gives
    $\|K\|^2=\|K^\dagger K\|=\|P_{\mathcal E}\|\leq1$.
    Also,
    \begin{align}
        \|K\psi\|^2
         & =\langle\psi,K^\dagger K\psi\rangle
        =\langle\psi,P_{\mathcal E}\psi\rangle
        =\|P_{\mathcal E}\psi\|^2.
    \end{align}
    Taking nonnegative square roots gives the norm equality. Orthogonal
    projections are contractive, so $\|K\psi\|\leq\|\psi\|$.
    Taking the supremum over unit vectors gives $\|K\|\leq1$.
\end{proof}

% Source: 05-frames.tex, eq:encoder-contraction, lines 61--64, at
% adc7f1241b42e322a6451854ab7e4b4c146bf78a.
% Auxiliary consequence; native verification of the new declaration is pending.
\begin{theorem}[Norm preservation on the cylinder span]
    \label{thm:peps_hole_encoder_norm_on_span}
    \lean{TNLean.PEPS.holeEncoder_norm_eq_of_mem}
    \uses{def:peps_hole_encoder}
    For the encoding $K$ of Definition~\ref{def:peps_hole_encoder}, every
    $\psi\in\mathcal E$ satisfies $\|K\psi\|=\|\psi\|$.
    This follows from the Gram identity in the
    \href{https://github.com/openai/math/blob/adc7f1241b42e322a6451854ab7e4b4c146bf78a/preprints/Polynomial-PEPS-approximation-of-gapped-square-grid-ground-states-September-24-2026/build/sections/05-frames.tex\#L61-L64}{one-hole encoding construction}.
\end{theorem}
\begin{proof}
    \uses{thm:peps_hole_encoder_norm,def:peps_coordinate_projector}
    Since $\psi\in\mathcal E$, its orthogonal projection satisfies
    $P_{\mathcal E}\psi=\psi$. Hence
    $\|K\psi\|=\|P_{\mathcal E}\psi\|=\|\psi\|$ by
    Theorem~\ref{thm:peps_hole_encoder_norm}.
\end{proof}

\begin{definition}[Identity encoding]
    \label{def:peps_identity_hole_encoder}
    \lean{
        TNLean.PEPS.identityHoleEncoder,
        TNLean.PEPS.identityHoleEncoder_mulVec
    }
    \leanok
    With a singleton tag set $\{\ast\}$, define
    $K_{\mathrm{id}}\psi=\ket{\ast}\otimes\psi$.
    This is the encoding used when an outer square has empty physical sample.
\end{definition}

\begin{theorem}[Identity Gram operator and norm]
    \label{thm:peps_identity_hole_encoder}
    \lean{
        TNLean.PEPS.identityHoleEncoder_gram,
        TNLean.PEPS.identityHoleEncoder_norm
    }
    \leanok
    \uses{def:peps_identity_hole_encoder}
    The identity encoding satisfies $K_{\mathrm{id}}^\dagger K_{\mathrm{id}}=I$.
    For every $\psi\in\mathcal H_V$,
    $\|K_{\mathrm{id}}\psi\|=\|\psi\|$.
\end{theorem}
\begin{proof}
    \leanok
    The sum over the singleton tag has one term, and the raw matrix is $I$.
    Thus the Gram matrix is $I^\dagger I=I$.
    Equality of Gram matrices gives equality of the vector norms.
\end{proof}

\begin{theorem}[Zero encodings]
    \label{thm:peps_hole_encoder_zero}
    \lean{
        TNLean.PEPS.holeEncoder_empty,
        TNLean.PEPS.holeEncoder_zero_spaces
    }
    \leanok
    \uses{def:peps_hole_encoder}
    If $n=0$ or every $S_j=0$, then $T$ is empty and $K=0$.
\end{theorem}
\begin{proof}
    \leanok
    \uses{thm:peps_hole_encoder_tag_dimension}
    An empty family has no branch index. If every $S_j=0$, the tag
    dimension bound gives $|T|\leq\sum_{j<n}\dim(0)=0$.
    Hence $T$ is empty. In either case there are no tag coordinates,
    so the matrix is zero.
\end{proof}

\begin{remark}[Empty families and empty regions]
    If $n=0$, the tag set $T$ is empty and $K=0$;
    also $\mathcal E=0$, in agreement with~(\ref{eq:peps_hole_encoder_gram}).
    More generally, a zero-dimensional projected image contributes no tags.
    For an empty region $R_j$, the inside configuration set is a singleton
    and $\mathcal H_{R_j}=\C$; the reset condition is automatic and the
    complement is all of $V$.
    The separate choice $K_{\mathrm{id}}$ has one tag and retains every
    raw configuration.
\end{remark}
